\section{Experimental Methodology}\label{sec:methodology}
\subsection{Research Questions and Experimental Units}
The extended assessment addresses four questions:
\begin{description}
\item[RQ1:] How do the symbolic/heuristic and LLM identification paths
compare in candidate quality, agreement, coverage, and analysis cost?
\item[RQ2:] What construction and functional outcomes are observed
after data handling and transformation, and which failure stages
account for unsuccessful cases?
\item[RQ3:] What are the observed outcomes and representation changes
for assembly wrapping in each assessed architecture?
\item[RQ4:] How do zero-shot and two-shot LLM transformation differ
in initial and final acceptance, functional outcomes, similarity,
diversity, time, and cost?
\end{description}

The input payload is the unit for comparisons across methods. A run
identifies one input and one configuration, while an attempt is a
generation or validation step within that run. Multiple variants of
one payload are repeated observations of the same input rather than
independent samples of the Windows software population. Scan engines
and repeated submissions of identical bytes are also not treated as
independent payloads.

\subsection{Corpora, Versions, and Configuration Matrix}
The conference corpus consists of 102 inputs, divided into 38 x64 and
64 x86 payloads~\cite{zone_secrypt_2026}. Extended campaigns are described
by their own manifests. A subset selected for poor historical outcomes
is identified as such; its aggregate result is not presented as a new
measurement on all 102 inputs.

Table~\ref{tab:campaign_manifest} reserves the information needed to
reconstruct each campaign. Distinct payload count, repeated runs,
constructed files, and unique output hashes remain separate counts.
The final record includes input hashes, engine version or snapshot
hash, architecture, prompt version, model configuration, and dependency
versions. For identification, the analysis limits and cache status are
also recorded.

\begin{table}[t]
\caption{Extended campaign manifest. Pending entries await the frozen
experimental datasets; the historical corpus is specified in the text.}
\label{tab:campaign_manifest}
\centering\small
\begin{tabularx}{\linewidth}{@{}Ylll@{}}
\toprule
Campaign & Inputs & Architectures & Replicates \\
\midrule
Code--data identification & \pending & x64/x86 & \pending \\
Data handling and native methods & \pending & x64/x86 & \pending \\
Assembly wrapping & \pending & x64 & \pending \\
Assembly wrapping port & \pending & x86 (development) & \pending \\
LLM transformation, 0/2 shot & \pending & \pending & \pending \\
\bottomrule
\end{tabularx}
\end{table}

Model aliases can change independently of the client. The identifier
sent in a request, the version returned by the service when available,
and any version label assigned locally are therefore distinguished.
A locally maintained tariff or alias table is not independent
confirmation of the model version used by the provider.

\subsection{Identification Quality and Reference Annotations}
Where reference annotations are available, data is treated as the
positive class. Let $G$ be the set of annotated data-byte offsets and
$P$ the set predicted by the method. The byte-level measures are
\begin{align}
\mathrm{Precision}_{D} &= \frac{|P\cap G|}{|P|}, &
\mathrm{Recall}_{D} &= \frac{|P\cap G|}{|G|},\\
F_{1,D} &= \frac{2\,\mathrm{Precision}_{D}\,\mathrm{Recall}_{D}}
 {\mathrm{Precision}_{D}+\mathrm{Recall}_{D}}.
\end{align}
Undefined cases, including inputs without reference data, are reported
explicitly rather than assigned a favorable score. Code bytes predicted
as data are counted as false positives and examined separately because
their treatment can invalidate execution.

Interval-level assessment complements byte-level scores. An exact
boundary match and partial overlap are distinguished; the chosen
overlap rule is stated with the results. Any reference bytes left
unresolved by annotation are excluded from binary classification
scores and their coverage is reported. The same treatment is used
for all methods.

Reference labels must be established independently of the evaluated
method and associated with the original input offsets. The assessment
retains raw candidates and post-filtering results separately. It can
therefore attribute an error to identification or to acceptance and
integration, rather than assigning every final outcome to the model.

If independent labels are unavailable, the comparison is reported as
agreement and a case study. Pairwise agreement between two methods,
or byte identity of their final outputs, does not establish accuracy.
Functionality after extraction is likewise reported as a downstream
check rather than used to construct reference labels.

\subsection{Build Success and Functional Outcomes}
Let $N$ be the number of planned cases, $B$ the number producing an
accepted binary, and $W$ the number passing the stated functional check.
The assessment reports
\begin{equation}
\mathrm{BuildRate}=B/N,\qquad
\mathrm{WorkingRate}_{\mathrm{all}}=W/N,\qquad
\mathrm{WorkingRate}_{\mathrm{built}}=W/B.
\end{equation}
All three counts accompany the percentages. An unbuilt case is not
removed from the first two denominators. An incomplete or unexecuted
case remains distinguishable from a completed failing test.

The functional oracle is specified by payload class. An observed
connection, an opened port, a dialog, and exception-free termination
are different observables with different evidential strength. The
harness version, architecture, observation window, and relevant
configuration are recorded. Original inputs are checked under the
same current harness to characterize its baseline behavior.

These observations do not prove complete semantic equivalence. A
connection can occur before later execution fails; clean termination
may not exercise the intended operation; and a dialog verifies only
the behavior observed by that test. Conclusions retain the scope of
the oracle actually used.

For LLM runs, the record additionally distinguishes initial acceptance,
acceptance within the retry budget, and the outcome of each functional
check. If execution is part of generation-time acceptance, a subsequent
pass rate on selected outputs is labeled accordingly. Exhausted runs
remain in the run-level denominator.

\subsection{Static-Scan Measurements}
For completed file $i$, let $m_i$ be the number of malicious verdicts
and $a_i$ the number of engines included by the stated eligibility rule.
The per-file malicious-verdict fraction and its macro average are
\begin{equation}
d_i=m_i/a_i,\qquad
\overline d=\frac{1}{F}\sum_{i=1}^{F}d_i.
\end{equation}
The file-level detected fraction is a separate measure,
\begin{equation}
q=\frac{1}{F}\sum_{i=1}^{F}\mathbf{1}[m_i>0].
\end{equation}
The term detection rate is not used to interchange these quantities.
The report specifies whether unsupported engines, failed analyses,
and timeouts enter $a_i$, and provides their counts. A pending
analysis has no measured verdict fraction and is not recorded as zero.

Files are submitted or retrieved by content identity. Repeated outputs
with one hash can reuse a scan, while that hash can still belong to
several experimental conditions. Deduplication of scanning is separated
from membership in result groups. The analysis states whether summaries
weight unique binaries, generation runs, or payloads equally.

Scan dates and service metadata are retained. Historical and current
campaigns may use different engine versions and submission contexts.
Differences across such campaigns are not attributed solely to the
transformation. Raw-binary verdicts are interpreted within their file
format and service setting, without extrapolation to deployed endpoint
detection.

\subsection{Similarity, Diversity, and Representation Size}
Every similarity score specifies its representation and reference.
The sequence-based validation score in the current LLM path uses
normalized instruction signatures and the matching-block ratio exposed
by \texttt{SequenceMatcher}~\cite{python_difflib}. If $M$ signature
elements occur in the selected matching blocks and the sequence lengths
are $u$ and $v$, the ratio is $2M/(u+v)$. This is not cosine similarity.

Normalization removes selected literal differences. In the NOP
condition the implementation also removes recognized no-op operations
before comparison. A score of one after that filtering indicates
agreement under the defined representation; it does not indicate
byte-identical files. The exact normalization version is retained with
the experiment.

For frequency-based comparison, let $A$ and $B$ be nonzero vectors
under a common token vocabulary. Cosine similarity is
\begin{equation}
\mathrm{CosSim}(A,B)=\frac{A\cdot B}{\|A\|\,\|B\|}.
\end{equation}
The report states whether tokens are mnemonics alone or include
operands, and whether the vectors describe an original/output pair
or pooled corpus text. These constructions are not combined in one
unqualified similarity column. With wrapping, the decoding view used
to obtain tokens is also stated.

Output diversity is measured independently of similarity to an input.
The basic exact measure reports distinct hashes among completed
outputs, both within a condition and across conditions. Repeated
generation does not imply distinct bytes. Size is recorded as bytes
and, for matched pairs, as $(L_{\mathrm{out}}-L_{\mathrm{in}})/L_{\mathrm{in}}$.
Any entropy statistic identifies whether its symbols are raw bytes or
assembly tokens; it remains descriptive rather than evidence of
functional preservation.

\subsection{Time, Cost, and Aggregation}
The timing record separates analysis, preparation, model calls, local
validation, and the stated total where these are available. Cached
analysis and reused outputs are distinguished from newly computed
runs. Costs include failed generation attempts as well as final
accepted artifacts. Their currency, tariff date, and estimation rule
are specified.

Results are stratified by architecture and input population before
aggregation. Repeated variants of the same input remain grouped when
estimating uncertainty. For small pilots, individual counts and
distributions are reported and broad population-level claims are
avoided. Any statistical comparison is accompanied by its unit of
analysis, uncertainty method, and treatment of missing observations.
\draftnote{Finalize the campaign manifests, reference annotations,
harness configuration, scan eligibility rule, model settings, and
aggregation policy from the frozen experiment records.}
